\documentclass[11pt,a4paper]{article}

\usepackage[margin=1in]{geometry}
\usepackage{amsmath,amssymb,amsthm}
\usepackage{microtype}
\usepackage{enumitem}
\usepackage{hyperref}
\usepackage{xcolor}
\usepackage{xspace}

\hypersetup{
  colorlinks=true,
  linkcolor=blue!50!black,
  citecolor=blue!50!black,
  urlcolor=blue!50!black,
}

\newtheorem{theorem}{Theorem}
\newtheorem{lemma}[theorem]{Lemma}
\newtheorem{proposition}[theorem]{Proposition}
\newtheorem{corollary}[theorem]{Corollary}
\theoremstyle{definition}
\newtheorem{definition}[theorem]{Definition}
\newtheorem{game}[theorem]{Game}
\theoremstyle{remark}
\newtheorem{remark}[theorem]{Remark}

\newcommand{\zap}{\textsf{ZAP}\xspace}
\newcommand{\zapmcp}{\textsf{ZAP-MCP}\xspace}
\newcommand{\zapaa}{\textsf{ZAP-A2A}\xspace}
\newcommand{\mcp}{\textsf{MCP}\xspace}
\newcommand{\aatwo}{\textsf{A2A}\xspace}
\newcommand{\rns}{\textsf{ZAP-RNS}\xspace}
\newcommand{\Adv}{\mathsf{Adv}}
\newcommand{\Exp}{\mathsf{Exp}}
\newcommand{\AEAD}{\textsf{AE}}
\newcommand{\KEM}{\textsf{KEM}}
\newcommand{\Sign}{\textsf{Sig}}
\newcommand{\negl}{\mathsf{negl}}

\title{\textbf{\zapmcp{} and \zapaa{}: Formal Security Proofs}}
\author{ZAP Protocol Authors\\
\href{https://zap-proto.io}{\texttt{zap-proto.io}}}
\date{2026}

\begin{document}
\maketitle

\begin{abstract}
We give formal security proofs for the constructions \zapmcp{} and
\zapaa{} described in the companion paper. We prove agent
non-repudiation under the EUF-CMA security of ML-DSA-65, agent
impersonation infeasibility under combined IND-CCA security of
ML-KEM-768 and EUF-CMA security of ML-DSA-65, and replay freshness
under the assumption that frame nonces are sampled with sufficient
entropy. The proofs proceed by a sequence-of-games argument in the
style of Shoup. Concrete advantage bounds are stated in terms of
adversarial query budgets; we instantiate them for the recommended
parameter set (ML-KEM-768, ML-DSA-65, AES-256-GCM, 96-bit nonces).
\end{abstract}

\section{Notation and Primitives}

\paragraph{Computational model.} All adversaries are probabilistic
polynomial-time (PPT) Turing machines. We work in the standard model
unless otherwise noted; ML-DSA-65 and ML-KEM-768 are NIST-standard
post-quantum primitives whose security claims are stated in
FIPS~204 and FIPS~203 respectively. We assume the recommended Category~3
parameter sets throughout.

\paragraph{Hash function.} $H : \{0,1\}^* \to \{0,1\}^{256}$ denotes
SHA3-256. Where the proofs require it, we model $H$ as a random oracle;
all use of $H$ in the construction is over collision-resistant input,
so the random-oracle assumption is conservative.

\paragraph{Signature scheme.} $\Sign = (\Sign.\mathsf{KGen},
\Sign.\mathsf{Sign}, \Sign.\mathsf{Vfy})$ is ML-DSA-65, modeled as
EUF-CMA secure with advantage
$\Adv^{\mathsf{eufcma}}_{\Sign}(\cdot)$. Concretely, NIST's claim is
that no PPT adversary achieves advantage greater than
$2^{-128}$ given $2^{64}$ signing queries and $2^{128}$ hash queries.

\paragraph{KEM.} $\KEM = (\KEM.\mathsf{KGen}, \KEM.\mathsf{Encaps},
\KEM.\mathsf{Decaps})$ is ML-KEM-768, modeled as IND-CCA secure with
advantage $\Adv^{\mathsf{indcca}}_{\KEM}(\cdot)$.

\paragraph{AEAD.} $\AEAD = (\AEAD.\mathsf{Seal}, \AEAD.\mathsf{Open})$ is
AES-256-GCM with 96-bit nonces, modeled as IND-CPA $+$ INT-CTXT (i.e.\
authenticated encryption) with advantage
$\Adv^{\mathsf{ae}}_{\AEAD}(\cdot)$.

\paragraph{Identity binding.} \rns{} record for an entity $X$ is a
tuple $R_X = (\mathit{kem\_pk}_X, \mathit{sig\_pk}_X, \mathit{caps}_X,
\sigma_{\mathsf{root}})$ where $\sigma_{\mathsf{root}}$ is a publisher
signature. We assume the publisher root is honest and its public key is
known to all participants out of band. (Compromise of the publisher
root is out of scope; mitigations are discussed at the end of the
companion paper.)

\section{Game-Based Definitions}

\subsection{Frame syntax}

A \zapmcp{}/\zapaa{} frame is the tuple
\[
F = (\mathit{magic}, v, s, n, c, t)
\]
where $\mathit{magic}$ is a 4-byte protocol tag, $v$ is the version
(1), $s = \mathit{sid}\,\Vert\,\mathit{seq}$ is the session and
sequence identifier, $n$ is the AEAD nonce, $c$ is the AEAD
ciphertext containing $\hat{\rho}\,\Vert\,\sigma$ where $\hat{\rho}$ is
the canonical CBOR payload and $\sigma$ is the per-frame signature,
and $t$ is the AEAD tag. We write $F_{A,T,\mathit{sid},\mathit{seq}}$
for a frame from $A$ to $T$ in session $\mathit{sid}$ at sequence
$\mathit{seq}$.

\subsection{Non-repudiation game}

\begin{game}[Non-repudiation $\Exp^{\mathsf{nonrep}}_{\zapmcp{}}$]
\label{game:nonrep}
The challenger generates honest agent keys
$(\mathit{sig\_sk}_A, \mathit{sig\_pk}_A) \leftarrow
\Sign.\mathsf{KGen}(1^\lambda)$ for some honest agent $A$, and
publishes $\mathit{sig\_pk}_A$. Adversary $\mathcal{A}$ is given access
to a signing oracle $\mathcal{O}_{\mathsf{Sign}}(m)$ that returns
$\sigma \leftarrow \Sign.\mathsf{Sign}(\mathit{sig\_sk}_A, m)$. Let
$\mathcal{Q}$ be the set of messages queried.

$\mathcal{A}$ outputs a tuple $(F^*, \hat{\rho}^*, \mathit{sid}^*,
\mathit{seq}^*, \sigma^*)$. The adversary wins if
\[
\Sign.\mathsf{Vfy}(\mathit{sig\_pk}_A,
H(\hat{\rho}^* \,\Vert\, \mathit{ctx}^*), \sigma^*) = 1
\quad \wedge \quad
\hat{\rho}^* \,\Vert\, \mathit{ctx}^* \notin \mathcal{Q}
\]
where $\mathit{ctx}^* = \mathtt{"zap-mcp/v1"} \,\Vert\, \mathit{sid}^*
\,\Vert\, \mathit{seq}^*$.
The advantage is
$\Adv^{\mathsf{nonrep}}_{\zapmcp{}}(\mathcal{A}) =
\Pr[\mathcal{A}\ \text{wins}]$.
\end{game}

\subsection{Impersonation game}

\begin{game}[Agent impersonation
$\Exp^{\mathsf{imp}}_{\zapmcp{}}$]\label{game:imp}
Setup as above; in addition, generate KEM key pair
$(\mathit{kem\_sk}_T, \mathit{kem\_pk}_T) \leftarrow
\KEM.\mathsf{KGen}(1^\lambda)$ for an honest tool $T$, and publish
$\mathit{kem\_pk}_T$ and $\mathit{sig\_pk}_A$. $\mathcal{A}$ may
\begin{itemize}[leftmargin=2em,itemsep=0.2em]
\item invoke $\mathcal{O}_{\mathsf{Sign}}(m)$ as before, restricted so
that no queried $m$ contains a session id matching the test session;
\item invoke $\mathcal{O}_{\mathsf{Decaps}}(c)$ for any $c$ except the
challenge ciphertext;
\item observe arbitrary protocol executions.
\end{itemize}
$\mathcal{A}$ wins if it completes a handshake transcript with $T$ that
$T$ accepts as authenticated to $A$, where $\mathcal{A}$ does not hold
$\mathit{sig\_sk}_A$.
\end{game}

\subsection{Replay game}

\begin{game}[Cross-session replay
$\Exp^{\mathsf{replay}}_{\zapmcp{}}$]\label{game:replay}
The challenger establishes two distinct honest sessions
$\mathit{sid}_1 \neq \mathit{sid}_2$ between $A$ and $T$. $\mathcal{A}$
observes the full traffic of $\mathit{sid}_1$. $\mathcal{A}$ then
attempts to inject a frame $F^* = F_{A,T,\mathit{sid}_1,\mathit{seq}}$
into $\mathit{sid}_2$ such that $T$ accepts it as a fresh frame in
$\mathit{sid}_2$. Advantage is the acceptance probability.
\end{game}

\section{Theorem 1: Non-Repudiation}

\begin{theorem}[Non-repudiation]\label{thm:nonrep}
Let $\mathcal{A}$ be a PPT adversary against $\Exp^{\mathsf{nonrep}}_{\zapmcp{}}$
making at most $q_s$ signing-oracle queries and $q_h$ hash queries.
There exists a PPT algorithm $\mathcal{B}$ against the EUF-CMA security
of ML-DSA-65 making at most $q_s$ signing queries, with
\[
\Adv^{\mathsf{nonrep}}_{\zapmcp{}}(\mathcal{A}) \;\leq\;
\Adv^{\mathsf{eufcma}}_{\Sign,\mathsf{ML\text{-}DSA\text{-}65}}(\mathcal{B})
+ \frac{q_h^2}{2^{257}}.
\]
The running time of $\mathcal{B}$ is within an additive constant of
that of $\mathcal{A}$.
\end{theorem}

\begin{proof}
We proceed by a sequence of games.

\textbf{Game 0.} The original game $\Exp^{\mathsf{nonrep}}_{\zapmcp{}}$. By definition,
$\Pr[\mathcal{A}\ \text{wins}\ \text{Game 0}] = \Adv^{\mathsf{nonrep}}_{\zapmcp{}}(\mathcal{A})$.

\textbf{Game 1.} Identical to Game 0 except that the challenger aborts
if the random oracle ever produces a collision: i.e., if at any point
two distinct queries $x \neq x'$ result in $H(x) = H(x')$. By the
birthday bound, $\Pr[\text{collision}] \leq q_h(q_h-1)/2^{257} \leq
q_h^2/2^{257}$. Therefore
\[
|\Pr[\mathcal{A}\ \text{wins}\ \text{Game 0}] - \Pr[\mathcal{A}\ \text{wins}\ \text{Game 1}]|
\leq \frac{q_h^2}{2^{257}}.
\]

\textbf{Game 2.} We construct $\mathcal{B}$. $\mathcal{B}$ runs the
EUF-CMA game against ML-DSA-65, receiving challenge public key
$\mathit{pk}^*$ and a signing oracle $\mathcal{O}^*_{\mathsf{Sign}}$.
$\mathcal{B}$ simulates Game 1 for $\mathcal{A}$ as follows:
\begin{itemize}[leftmargin=2em,itemsep=0.2em]
\item Set $\mathit{sig\_pk}_A := \mathit{pk}^*$ and publish.
\item On $\mathcal{O}_{\mathsf{Sign}}(m)$: query
$\mathcal{O}^*_{\mathsf{Sign}}(m)$, return the result, and add $m$ to
$\mathcal{Q}$.
\item On a hash query, simulate the random oracle by lazy table
maintenance with the no-collision invariant from Game 1.
\end{itemize}
When $\mathcal{A}$ outputs $(F^*, \hat{\rho}^*, \mathit{sid}^*,
\mathit{seq}^*, \sigma^*)$, $\mathcal{B}$ computes
$m^* := H(\hat{\rho}^* \,\Vert\, \mathit{ctx}^*)$ and outputs
$(m^*, \sigma^*)$ as its forgery.

\textbf{Reduction analysis.} If $\mathcal{A}$ wins Game 1, then by the
winning condition $\sigma^*$ is a valid signature on $m^*$, and
$\hat{\rho}^* \,\Vert\, \mathit{ctx}^* \notin \mathcal{Q}$. Since the
random oracle is collision-free in Game 1, $m^*$ is therefore distinct
from every $H(m)$ for $m \in \mathcal{Q}$, so $\mathcal{B}$ never
queried $m^*$ to $\mathcal{O}^*_{\mathsf{Sign}}$. Hence
$(m^*, \sigma^*)$ is a valid EUF-CMA forgery against ML-DSA-65, and
\[
\Pr[\mathcal{B}\ \text{wins}] = \Pr[\mathcal{A}\ \text{wins}\ \text{Game 1}].
\]
Combining the two bounds:
\[
\Adv^{\mathsf{nonrep}}_{\zapmcp{}}(\mathcal{A}) \leq
\Adv^{\mathsf{eufcma}}_{\Sign}(\mathcal{B}) + \frac{q_h^2}{2^{257}}.
\]
\qed
\end{proof}

\begin{corollary}[Concrete instantiation]\label{cor:nonrep}
With NIST's claim that
$\Adv^{\mathsf{eufcma}}_{\mathsf{ML\text{-}DSA\text{-}65}}(\mathcal{B}) \leq 2^{-128}$
for any PPT $\mathcal{B}$ with $q_s \leq 2^{64}$, and with $q_h \leq
2^{128}$,
\[
\Adv^{\mathsf{nonrep}}_{\zapmcp{}}(\mathcal{A}) \leq 2^{-128} + 2^{-1} \cdot 2^{-1}
= 2^{-128} + \tfrac{1}{4}.
\]
The dominant term for any realistic $q_h$ is $2^{-128}$. For
deployment-realistic $q_h \leq 2^{80}$, the bound tightens to
$2^{-128} + 2^{-97}$, i.e.\ negligibly above $2^{-97}$.
\end{corollary}

\section{Theorem 2: Agent Impersonation Infeasibility}

\begin{theorem}[Agent impersonation]\label{thm:imp}
Let $\mathcal{A}$ be a PPT adversary against
$\Exp^{\mathsf{imp}}_{\zapmcp{}}$ making at most $q_d$ decapsulation
queries, $q_s$ signing queries, and $q_h$ hash queries. There exist
PPT algorithms $\mathcal{B}_1$ against IND-CCA of ML-KEM-768 and
$\mathcal{B}_2$ against EUF-CMA of ML-DSA-65 with
\[
\Adv^{\mathsf{imp}}_{\zapmcp{}}(\mathcal{A}) \;\leq\;
\Adv^{\mathsf{indcca}}_{\KEM,\mathsf{ML\text{-}KEM\text{-}768}}(\mathcal{B}_1)
+ \Adv^{\mathsf{eufcma}}_{\Sign,\mathsf{ML\text{-}DSA\text{-}65}}(\mathcal{B}_2)
+ \frac{q_h^2}{2^{257}}.
\]
\end{theorem}

\begin{proof}
The handshake transcript that $T$ accepts contains, by construction:
\begin{enumerate}[leftmargin=2em,itemsep=0.2em]
\item $\mathcal{A}$'s presented identity claim, naming $A$ and
including a signature
$\sigma_h = \Sign.\mathsf{Sign}(\mathit{sig\_sk}_A, H(\mathit{th}))$
over the running hash $\mathit{th}$ of the transcript;
\item the KEM ciphertext $c_T = \KEM.\mathsf{Encaps}(\mathit{kem\_pk}_T)$
contributing to the session secret;
\item key confirmation MAC over the handshake-derived session key.
\end{enumerate}
For $T$ to accept, $\sigma_h$ must verify under
$\mathit{sig\_pk}_A$. Two cases arise.

\textbf{Case A: $\mathit{th}$ has been signed by $A$ in some earlier
session.} By the construction of the transcript hash, $\mathit{th}$
incorporates $T$'s nonce $r_T$, the session identifier
$\mathit{sid}$, and a freshly-sampled handshake-time random from $A$.
Two transcripts $\mathit{th} = \mathit{th}'$ would imply collision in
$H$, occurring with probability $\leq q_h^2 / 2^{257}$.

\textbf{Case B: $\mathit{th}$ has never been signed by $A$.} Then
$\sigma_h$ is a forgery of ML-DSA-65 on a fresh message under
$\mathit{sig\_pk}_A$. We construct $\mathcal{B}_2$ as in
Theorem~\ref{thm:nonrep}, embedding the EUF-CMA challenge public key
as $\mathit{sig\_pk}_A$ and answering $\mathcal{A}$'s signing queries
through the EUF-CMA oracle. We have
\[
\Pr[\text{Case B}] \leq \Adv^{\mathsf{eufcma}}_{\Sign}(\mathcal{B}_2).
\]

\textbf{Reducing key confirmation to KEM IND-CCA.} If $\mathcal{A}$
manages to compute the correct key confirmation MAC without holding
the KEM secret, then either (i) it has guessed the shared secret, or
(ii) it has obtained it through the decapsulation oracle. Case (i) is
information-theoretically negligible ($2^{-256}$ for a 256-bit
shared secret). For Case (ii), we construct $\mathcal{B}_1$ that
embeds the IND-CCA challenge public key as $\mathit{kem\_pk}_T$ and
answers decapsulation queries through its own oracle (with the usual
restriction excluding the challenge ciphertext). $\mathcal{B}_1$
distinguishes the real shared secret from a uniform random by
observing whether $\mathcal{A}$ produces the matching MAC. We have
\[
\Pr[\text{key confirmation forged via KEM}] \leq
\Adv^{\mathsf{indcca}}_{\KEM}(\mathcal{B}_1).
\]

\textbf{Union bound.} Combining,
\[
\Adv^{\mathsf{imp}}_{\zapmcp{}}(\mathcal{A}) \leq
\Adv^{\mathsf{indcca}}_{\KEM}(\mathcal{B}_1) +
\Adv^{\mathsf{eufcma}}_{\Sign}(\mathcal{B}_2) +
\frac{q_h^2}{2^{257}}.
\]
\qed
\end{proof}

\begin{corollary}[Concrete impersonation bound]
For NIST Category~3 parameters,
$\Adv^{\mathsf{imp}}_{\zapmcp{}}(\mathcal{A}) \leq 2^{-128} + 2^{-128} +
q_h^2/2^{257} \approx 2^{-127}$ for $q_h \leq 2^{64}$.
\end{corollary}

\section{Lemma: Replay Freshness}

\begin{lemma}[Replay freshness]\label{lem:replay}
Let $\nu$ be the AEAD nonce length in bits, and let $Q$ be the maximum
per-session frame count. For two distinct sessions $\mathit{sid}_1
\neq \mathit{sid}_2$, the probability that any single frame from
$\mathit{sid}_1$ is accepted as fresh in $\mathit{sid}_2$ is at most
\[
\Pr[\mathit{replay\ accepted}] \leq
\Adv^{\mathsf{ae}}_{\AEAD}(\mathcal{B}) + \frac{Q(Q-1)}{2^{\nu+1}}.
\]
\end{lemma}

\begin{proof}
The AEAD nonce in \zap{} is derived as
$\mathit{nonce}(\mathit{sid}, \mathit{seq}) = \mathsf{HKDF}(K_n,
\mathit{sid} \,\Vert\, \mathit{seq})$ where $K_n$ is a session-bound
HKDF key derived from the handshake-shared secret. Two distinct
sessions yield independently derived $K_n$ values; thus an injected
frame from $\mathit{sid}_1$ presents a nonce that is, in the
$\mathit{sid}_2$ session's nonce-space, indistinguishable from
uniformly random unless the AEAD construction is broken.

The injected frame's AEAD tag was computed under $K_{AB}^{(1)}$, the
$\mathit{sid}_1$ session key. The recipient verifies under $K_{AB}^{(2)}$,
the $\mathit{sid}_2$ session key. By the INT-CTXT property of AEAD,
the probability of a tag verifying under a different key is at most
$\Adv^{\mathsf{ae}}_{\AEAD}(\mathcal{B})$, where $\mathcal{B}$ is a
straightforward reduction.

The remaining term, $Q(Q-1)/2^{\nu+1}$, bounds the probability of a
nonce collision \emph{within} session $\mathit{sid}_2$ if the
injection's nonce happens to clash with a pre-existing legitimate
frame's nonce. By the birthday bound on $\nu$-bit nonces with $Q$
samples, this is at most $Q(Q-1)/2^{\nu+1}$.
\qed
\end{proof}

\begin{corollary}
With $\nu = 96$ and $Q \leq 2^{40}$,
$\Pr[\mathit{replay\ accepted}] \leq 2^{-128} + 2^{-17} \cdot 2^{-40} =
2^{-128} + 2^{-57}$, dominated by $2^{-57}$, well below operationally
relevant thresholds. Operators desiring tighter bounds may rotate
sessions before $Q = 2^{32}$, yielding $\leq 2^{-32} \cdot 2^{-32} =
2^{-65}$ replay probability.
\end{corollary}

\section{Composite Property: Provenance Chain Integrity for \zapaa{}}

We extend Theorem~\ref{thm:nonrep} to multi-hop \aatwo{} task chains.

\begin{theorem}[Provenance chain integrity]\label{thm:prov}
Let $\tau$ be an \aatwo{} task delegated through agents
$A_1 \to A_2 \to \dots \to A_k$, with each delegation producing a
signed frame as described in the companion paper. Let $\mathcal{A}$
be a PPT adversary that produces a forged delegation chain
$(F_1^*, F_2^*, \dots, F_k^*)$ with at least one $F_i^*$ purporting to
be from honest $A_i$ but not actually signed by $A_i$. Then
\[
\Adv^{\mathsf{prov}}_{\zapaa{}}(\mathcal{A}) \leq
k \cdot \left(
\Adv^{\mathsf{eufcma}}_{\Sign}(\mathcal{B}_i) + \frac{q_h^2}{2^{257}}
\right).
\]
\end{theorem}

\begin{proof}
By a hybrid argument over the $k$ positions in the chain. For each
position $i$, an adversary that succeeds with non-negligible advantage
on position $i$ can be reduced to an EUF-CMA adversary on $A_i$'s key
by Theorem~\ref{thm:nonrep}. The factor $k$ arises from the union bound
over positions.
\qed
\end{proof}

\section{Composite Property: Capability-Bound Method Restriction}

\begin{proposition}[Capability binding]\label{prop:caps}
Suppose tool $T$'s \rns{} record advertises capability set
$\mathit{caps}_T \subseteq \mathcal{M}$, where $\mathcal{M}$ is the
universe of \mcp{} method names. The handshake transcript hash binds
$H(\mathit{caps}_T)$. If a frame in the resulting session invokes a
method $m \notin \mathit{caps}_T$, an honest tool $T$ rejects without
processing.
\end{proposition}

\begin{proof}[Proof sketch]
Direct from construction: the handshake binding includes a commitment
to $\mathit{caps}_T$ verifiable by both endpoints, and the framed
method invocation is checked against $\mathit{caps}_T$ before
dispatch. The only way to bypass this is to forge a transcript
admitting a different $\mathit{caps}_T$, which requires either
forging the publisher root signature (out of scope) or a hash
collision in $H$ ($\leq q_h^2 / 2^{257}$).
\qed
\end{proof}

\section{Forward Secrecy of Past Sessions}

\begin{proposition}[Forward secrecy]\label{prop:fs}
Suppose at time $t$ all long-term keys
$\mathit{sig\_sk}_A, \mathit{kem\_sk}_T$ are exposed to $\mathcal{A}$.
For any session that completed at time $t' < t$ with non-overlapping
ephemeral handshake material, $\mathcal{A}$'s advantage at recovering
the session payload is bounded by
\[
\Adv^{\mathsf{fs}}_{\zapmcp{}}(\mathcal{A}) \leq
\Adv^{\mathsf{indcca}}_{\KEM,\mathsf{ephemeral}}(\mathcal{B})
+ \Adv^{\mathsf{ae}}_{\AEAD}(\mathcal{B}').
\]
\end{proposition}

\begin{proof}[Proof sketch]
\zap{} handshake derives the session secret from a hybrid combination
of static-static, static-ephemeral, and ephemeral-ephemeral KEM
encapsulations. The ephemeral-ephemeral component, by construction
discarded after the handshake, is independent of the long-term keys.
Compromise of long-term keys at time $t$ does not recover ephemerals
discarded at time $t' < t$. Reduction is to ephemeral KEM IND-CCA.
\qed
\end{proof}

\section{Summary Table}

\begin{table}[h]
\centering
\small
\renewcommand{\arraystretch}{1.2}
\begin{tabular}{|l|c|}
\hline
\textbf{Property} & \textbf{Bound} \\
\hline
Non-repudiation (Thm.~\ref{thm:nonrep})       & $\Adv^{\mathsf{eufcma}}_{\Sign} + q_h^2/2^{257}$ \\
Impersonation (Thm.~\ref{thm:imp})            & $\Adv^{\mathsf{indcca}}_{\KEM} + \Adv^{\mathsf{eufcma}}_{\Sign} + q_h^2/2^{257}$ \\
Replay (Lem.~\ref{lem:replay})                & $\Adv^{\mathsf{ae}}_{\AEAD} + Q(Q-1)/2^{\nu+1}$ \\
Provenance chain (Thm.~\ref{thm:prov})        & $k \cdot \left(\Adv^{\mathsf{eufcma}}_{\Sign} + q_h^2/2^{257}\right)$ \\
Capability binding (Prop.~\ref{prop:caps})    & $q_h^2/2^{257}$ \\
Forward secrecy (Prop.~\ref{prop:fs})         & $\Adv^{\mathsf{indcca}}_{\KEM,\mathsf{ephemeral}} + \Adv^{\mathsf{ae}}_{\AEAD}$ \\
\hline
\end{tabular}
\end{table}

\section{Discussion}

\subsection{Tightness}

Theorem~\ref{thm:nonrep}'s reduction is essentially tight: any
non-repudiation forgery in \zapmcp{} translates directly to an
ML-DSA-65 forgery, with only the random-oracle collision term
separating the two probabilities. We have therefore reduced the
question of agent message non-repudiation to the question of
ML-DSA-65 EUF-CMA security, for which NIST's standardisation body
provides the conservative claim of $\geq 192$ bits of post-quantum
security.

\subsection{Standard model vs random oracle}

The random oracle assumption appears only in the bookkeeping of the
hash function used to bind context strings. Standard-model
construction is straightforward by replacing $H$ with a
collision-resistant hash plus a domain separation tag, at the cost
of marginally larger advantage terms. We have chosen to state proofs
in the simpler model.

\subsection{Quantum oracle access}

The above proofs are stated in the classical-query model.
Post-quantum security claims for ML-DSA-65 and ML-KEM-768 are stated
under quantum-oracle access in the relevant FIPS standards. The
reductions above are tight under classical-query reductions and
remain tight up to a constant factor under known-quantum lifts
(we elide the technical details for brevity).

\subsection{What is not proven}

We do not prove anything about the semantics of the JSON-RPC payloads
themselves, nor about the absence of side-channel leakage from agent
implementations of ML-DSA-65 signing. Both belong to layers above and
below the protocol respectively. Additionally, model-layer
prompt-injection attacks are out of scope; \zapmcp{} provides
cryptographic attribution but not semantic safety.

\section{Conclusion}

We have given concrete-advantage bounds for the principal security
properties of \zapmcp{} and \zapaa{}: non-repudiation under
ML-DSA-65 EUF-CMA, agent impersonation infeasibility under
ML-KEM-768 IND-CCA combined with ML-DSA-65 EUF-CMA, replay freshness
under AEAD INT-CTXT, multi-hop provenance via a hybrid argument,
capability binding via transcript hashing, and forward secrecy of
past sessions via ephemeral KEM material. All bounds are
parameterised by adversarial query budgets and instantiate to
$\geq 128$-bit post-quantum security for the recommended
parameter set.

\end{document}
